%% ma: L12-B06-0037
%% lop: 12
%% chuong: 2
%% bai: 6
%% dang: 12.6.14
%% loai: TLN
%% muc: TH
%% dap_an: "70"
%% nguon: "LaTeX giáo viên gửi 21/10/2026 (M05, Vectơ và hệ trục tọa độ trong không gian, 107 câu), Phần III Câu 6"
%% kien_thuc: "Bài 6 (tổng vectơ, bình phương vô hướng)"
%% trang_thai: cho-duyet
%% ly_do: "Dạng mới đề xuất 12.6.14"
%% ghi_chu: "Đề gốc viết 'lực cản F4', đã viết 'lực F4' (không đổi ý)"
\begin{ex}
Một thiết bị bay flycam chịu tác động bởi ba lực gió $\overrightarrow{F_1}$, $\overrightarrow{F_2}$, $\overrightarrow{F_3}$ có phương đôi một vuông góc với nhau và có độ lớn lần lượt là $20$ N, $30$ N, $60$ N. Để thiết bị đứng yên cân bằng, động cơ cần sinh ra một lực $\overrightarrow{F_4}$ ngược hướng với hợp lực của gió. Tính độ lớn của lực $\overrightarrow{F_4}$ (đơn vị: N).
\shortans{70}
\loigiai{$|\overrightarrow{F_1}+\overrightarrow{F_2}+\overrightarrow{F_3}|=\sqrt{20^2+30^2+60^2}=70$ nên $|\overrightarrow{F_4}|=70$ N.}
\end{ex}
